\section*{Chapter \ref{ch:b3:difference-in-differences} --- Difference-in-Differences}

\begin{solution}{exo:b3:difference-in-differences:1}
\begin{enumerate}
\item Pennsylvania: $21.17 - 23.33 = -2.16$. New Jersey: $21.03 - 20.44 = 0.59$.
Difference in differences: $0.59 - (-2.16) = 2.75$ workers per restaurant (the
paper's 2.76 comes from unrounded means). % figdata: ck_stats
\item $0.59/20.44 = 0.0289$: employment rose by 2.9~\%.
\item That without the rise, New Jersey's employment per restaurant would have
fallen by as much as Pennsylvania's did, from 20.44 to 18.28.
\end{enumerate}
\end{solution}

\begin{solution}{exo:b3:difference-in-differences:2}
\begin{enumerate}
\item First region: $21 - 20 = 1$. Second: $26.5 - 24 = 2.5$. Estimate: $2.5 - 1 =
\dal{1.5}$ an hour.
\item Parallel trends: without the cut, wages in the second region would have
risen by 1 dalen like the first's. Check that the two regions' wages moved in
parallel in the years before 2024.
\end{enumerate}
\end{solution}

\begin{solution}{exo:b3:difference-in-differences:3}
At $k = 2$ the estimate is 0.69 with an interval of 0.56 to 0.82, covering the
true 0.7. % figdata: event_study
Leads rising steadily from $-0.4$ to $-0.1$ would show the treated units
already gaining on the others before adoption, at about 0.1 a period. Parallel
trends would be rejected, and the lags would mix that trend with the effect.
\end{solution}

\begin{solution}{exo:b3:difference-in-differences:4}
Contradicting: a gap between the groups that widens or narrows steadily before
the treatment (a differential trend); a dip in the treated group just before
treatment (Ashenfelter's dip). Not contradicting: a constant gap in levels, with
both groups moving up and down together.
\end{solution}

\begin{solution}{exo:b3:difference-in-differences:5}
\begin{enumerate}
\item Levels: $20 - 10 = 10$. Logs: $\ln(120/100) - \ln(60/50) = \ln 1.2 -
\ln 1.2 = 0$.
\item Parallel trends is about A's counterfactual: in levels, A without the
treatment would have gained 10, like B; in logs, it would have grown 20~\% like B,
that is by 20. Starting from different levels, a gain of 10 and a growth of 20~\%
are different numbers for A, so at most one can hold. Choose the scale on which the groups moved in
parallel before the treatment, and on which the outcome's changes are
comparable across its levels: for sales, percentage changes usually.
\end{enumerate}
\end{solution}

\begin{solution}{exo:b3:difference-in-differences:6}
The tested solution is \code{ex06_placebo_1974.py}, in the chapter's
\code{exercises/} folder.
\begin{enumerate}
\item Trainees' change $-\usd{564}$; population survey $-\usd{366}$,
placebo estimate $-198$ (s.e.\ 280); the second survey $-\usd{365}$, placebo
$-198$ (s.e.\ 317). % output: ex06_placebo_1974
(From the printed changes the second difference is $-199$; the script works with
unrounded means.)
\item Both estimates are small relative to their standard errors: the
trainees' earnings did not fall significantly faster than the surveys' before the
programme. That is consistent with parallel trends from 1974 to 1975. It cannot
show they held from 1975 to 1978, and it does not rule out a dip that began
in 1975 itself, after the date the data measure.
\end{enumerate}
\end{solution}

\begin{solution}{exo:b3:difference-in-differences:7}
Upward. The before year is abnormally low, so part of the rise afterwards is a
recovery that would have happened anyway, and the comparison group, without a
dip, has no such recovery to subtract. Using a before year earlier than the dip
compares from a normal level and removes that part of the bias.
\end{solution}

\begin{solution}{exo:b3:difference-in-differences:8}
\begin{enumerate}
\item The dummy is E: $(0, 1, 1)$, L: $(0, 0, 1)$. Group means: E $2/3$, L
$1/3$; period means $0, 1/2, 1$; overall mean $1/2$. Subtracting the group and
period means and adding the overall mean gives E: $(-1/6, 1/3, -1/6)$ and L:
$(1/6, -1/3, 1/6)$.
\item The sum of squares is $4 \times 1/36 + 2 \times 1/9 = 1/3$. Unit and period
effects in $Y$ are orthogonal to the demeaned dummy, so only the treatment
effects count: $\sum \tilde D Y = \frac13 \tau - \frac16(\tau + \gamma) + \frac16
\tau = \frac{\tau}{3} - \frac{\gamma}{6}$. Dividing by $1/3$ gives $\tau -
\gamma/2$.
\item With $\tau = 1$ and $\gamma = 3$ the coefficient is $-0.5$, although every
effect is positive. The cell weights are 1 on (E, period 2), $-1/2$ on (E,
period 3) and $1/2$ on (L, period 3); the negative weight is on the early
group's later, larger effect, which serves as a comparison for L.
\end{enumerate}
\end{solution}

\begin{solution}{exo:b3:difference-in-differences:9}
Placebo group: provinces far from the western border, split at random into
``treated'' and ``untreated'', or an eastern province given a fictitious tax
date; an effect there would show a shock specific to some provinces. Placebo
outcome: sales of goods the tax did not touch, such as groceries; a fall there
would point to a regional downturn rather than the tax. Placebo date: pretend
the tax came two years earlier and use only the years before the true date; an
effect would reveal different pre-existing trends.
\end{solution}

\begin{solution}{exo:b3:difference-in-differences:10}
The tested solution is \code{ex10_anticipation.py}, in the chapter's
\code{exercises/} folder.
\begin{enumerate}
\item With $k = -1$ as reference: $-0.36$, $-0.31$, $0.16$, $0.28$ and $0.39$ at
$k = -3, -2, 0, 1, 2$. % output: ex10_anticipation
\item With $k = -2$ as reference: $-0.05$ at $k = -3$, $0.31$ at $k = -1$, and
$0.47$, $0.59$, $0.70$ at $k = 0, 1, 2$, close to the truth (0.3, then 0.5, 0.6,
0.7). In the first regression every coefficient is measured against a period
already raised by 0.3, so all shift down by about 0.3: the leads look like a
pre-trend and the effects are understated. The reference period must be one that
the treatment cannot have reached.
\end{enumerate}
\end{solution}

\begin{solution}{exo:b3:difference-in-differences:11}
\begin{enumerate}
\item Solve $2w_A + 6w_B + 4w_C = 4$, $4w_A + 6w_B + 10w_C = 6$, $w_A + w_B + w_C =
1$. Substituting $w_C = 1 - w_A - w_B$, the first gives $w_A = w_B$, the second
$-6w_A - 4w_B = -4$, so $w_A = w_B = 0.4$ and $w_C = 0.2$, all non-negative.
Check: $0.8 + 2.4 + 0.8 = 4$ and $1.6 + 2.4 + 2 = 6$.
\item Synthetic: $0.4 \times 3 + 0.4 \times 7 + 0.2 \times 9 = 5.8$. Effect: $5 -
5.8 = -0.8$.
\item Negative weights, or weights above one, would let the synthetic unit lie
outside the range of the donors: an extrapolation that can fit the past exactly
and predict the future badly, and that no reader could interpret as a
comparison with real places.
\end{enumerate}
\end{solution}

\begin{solution}{exo:b3:difference-in-differences:12}
\begin{enumerate}
\item With $\Delta$ the after-minus-before change of a group's mean, the
estimator is $[\Delta(\text{old, state}) - \Delta(\text{young, state})] -
[\Delta(\text{old, others}) - \Delta(\text{young, others})]$.
\item It needs the gap between the old's and the young's untreated trends to be
the same in the state as elsewhere. That allows any shock to the whole state
(hitting old and young alike) and any nationwide shock specific to the old,
both of which bias the simple comparison of the old; it fails only for a shock
that hits the old in that state alone.
\end{enumerate}
\end{solution}

\begin{solution}{pb:b3:difference-in-differences:border}
\textbf{1.} West: $5.0773 - 5.1633 = -0.0860$. North: $5.1788 - 5.1871 =
-0.0083$. Difference in differences: $-0.0860 - (-0.0083) = -0.0777$.
% figdata: kestria_means

\textbf{2.} $\eu^{-0.0777} - 1 = -0.0748$: sales 7.5~\% lower.

\textbf{3.} $-0.0777 \pm 1.96 \times 0.0225 = -0.0777 \pm 0.0441$: from $-0.1218$
to $-0.0336$.

\textbf{4.} Nine: three western and six northern provinces. Cluster-robust
standard errors are justified as the number of clusters grows and tend to be too
small with few; with three treated clusters, the interval should be checked by
methods that resample or permute whole provinces.

\textbf{5.} No. Parallel trends allow any difference in levels; only a
difference in changes would bias the estimate.

\textbf{6.} For instance, a large shopping centre opening on the Brenland side
of the border in 2019 would draw western shoppers whatever the tax, and make
the estimate more negative. A downturn in the west's main industry would do the
same.

\textbf{7.} Dividing each lead by its standard error, from $k = -5$ to $-2$:
$-0.0088/0.0211 = -0.42$, $-0.0004/0.0221 = -0.02$, $-0.0036/0.0173 = -0.21$
and $-0.0163/0.0217 = -0.75$. None is near 2: the west
and north moved together before the tax. % figdata: kestria_event

\textbf{8.} The panel is balanced, every western shop is treated in the same
year, and each event-time dummy picks one year. With shop and year effects, the
coefficient on one year's dummy is the difference in differences between that
year and the reference year; $k = 1$ is 2020 against 2018, the comparison of
Part~I.

\textbf{9.} $-0.030$ in 2019, $-0.078$ in 2020, $-0.112$ in 2021, then $-0.146$,
$-0.157$ and $-0.144$: the effect builds for about three years and then levels
off near $-0.15$. Shoppers learn of the price gap, try the trip, and change
habits; shops across the border adapt their opening hours and stock.

\textbf{10.} Each coefficient compares one year with 2018. Persistent province
shocks drift apart as the two years grow further apart, so the comparison
carries more noise; and with nine clusters the standard errors are themselves
imprecise.

\textbf{11.} Sales in the west would be high in 2018, the reference year. Every
coefficient would be measured from an inflated base: the leads would turn
negative, as if the west had been rising before the tax, and the lags would
overstate the fall.

\textbf{12.} The east is treated from 2021. In a regression whose dummies mark
only western event times, eastern shops would count as untreated in 2021--2024,
and their fall would be read as a fall of the comparison group.

\textbf{13.} The west against the north; the east against the north; the west
against the east before 2021, while the east is untreated; and the east
against the west from 2021, while the west is already treated.

\textbf{14.} The last. From 2020 to 2021 the west's own effect deepens, from
$-0.08$ to $-0.12$ in the design; subtracting the west's change from the east's
subtracts that deepening, which makes the east's effect look smaller.

\textbf{15.} $(-0.0907) - (-0.106) = 0.0153$, and $0.0153/0.106 = 0.144$: \omterm{def:b3:panel-data:twfe}{two-way
fixed effects} understates the effect by 14~\%. % figdata: kestria_estimates

\textbf{16.} It compares each treated province only with the never-treated north,
and from its own last untreated year: no already-treated province ever serves
as a comparison.

\textbf{17.} 0.0225 against 0.0102; the two-way regression uses all years and
provinces. A smaller standard error around a biased number is no reason to
prefer it: the two-way estimate's interval, $-0.0907 \pm 0.0200$, from
$-0.1107$ to $-0.0707$, covers the truth only near its edge.

\textbf{18.} No. If each province's effect were constant over time, the
already-treated comparison group's change would contain no change in its
effect, and the forbidden comparisons would be harmless; differences across
provinces would affect only how the effects are weighted.

\textbf{19.} $0.5184 + 0.2863 + 0.1231 + 0.0722 = 1.0000$; four donors, N1--N4,
receive weight. % figdata: kestria_synth

\textbf{20.} The squares of the five gaps sum to $0.00124087$; divided by 5 and
square-rooted, $0.0158$. % figdata: kestria_synth_gap

\textbf{21.} $(-0.1527 - 0.1479 - 0.1975 - 0.1846)/4 = -0.1707$. The west's lags
for 2021--2024 ($k = 2$ to 5) average $-0.5589/4 = -0.1397$: W1's gap is larger.

\textbf{22.} W1 is the province where the trip to Brenland is shortest, so more
of its shoppers can switch.

\textbf{23.} The squares of the six gaps for 2019--2024 sum to $0.13078737$;
divided by 6 and square-rooted, $0.1476$. Ratio: $0.1476/0.0158 = 9.3$.

\textbf{24.} With W1 and six donors there are seven ratios; the smallest
$p$-value is $1/7 = 0.14$. A rank of 1 is the strongest evidence the test can
give, but it cannot reach the 5~\% level.

\textbf{25.} The east is taxed from 2021, so its later path contains the tax
effect. A synthetic W1 built partly from eastern provinces would fall with
them after 2021 and understate W1's effect.

\textbf{26.} The effect of the tax on western shops' sales: $-0.078$ by the
two-by-two comparison of 2020 with 2018, when the effect was still building;
$-0.091$ by \omterm{def:b3:panel-data:twfe}{two-way fixed effects}, biased toward zero by the roll-out to the
east; $-0.106$ by clean comparisons with the north, the number to report, with
the event-study path beside it; and, for the border province W1 alone,
$-0.171$ a year over 2021--2024 by \omterm{def:b3:difference-in-differences:synthetic-control}{synthetic control}, with a \omterm{def:b3:difference-in-differences:placebo}{placebo test} that
cannot be significant at 5~\% with six donors.
\end{solution}
